\chapter{SEAL 子图模式提取}\label{chap:seal}

第~\ref{chap:hgt}~章给出的异构图变换器在整图层面学习节点表征，再以拼接式评分头输出婚配 logit。然而对历史学家而言，单一标量得分难以提供"为何此对会被推荐"的局部依据。为此，本章在主管线之外另起一支独立的子图链路预测分支，把每一候选对周围的局部拓扑显式提取出来，并以可命名的微结构图样形式向下游多智能体协商提供结构指纹。本章依次给出任务定式、双半径节点标注、十三类历史微结构图样目录、关系图卷积评分模型，以及训练—评估实证。

\section{任务定式}\label{sec:seal-formulation}

给定候选对 $(m,w,\tau)$，其中 $m$ 为候选男方、$w$ 为候选女方、$\tau$ 为目标婚姻年份，子图链路预测（Subgraph Embedding And Labeling，SEAL）\citep{zhang2018seal} 把"$m$ 与 $w$ 在 $\tau$ 年是否结为婚配"的链路预测问题重新表述为一个以二元组为中心的封闭子图二分类问题。具体而言，从因果可见图 $\mathcal{G}^{\mathrm{vis}}_\tau$（即所有 $\mathit{edge\_time}<\tau$ 的边构成的快照、且母系边 $r_\mathrm{ms}, r_\mathrm{md}$ 已删）出发，沿无向投影抽取以 $m$ 与 $w$ 为双中心的 $h$ 跳封闭子图，再依次施加六条反泄漏屏蔽以删除可能直接揭示婚姻标签的边，得到掩码子图 $\widetilde{\mathcal{G}}(m,w,\tau)$。模型的任务是给该子图打一个标量结构得分。

值得注意的是，子图链路预测与异构图变换器的婚姻评分构成互补关系。前者把每对候选的判据约束在直径不超过 $h$ 跳的局部邻域内，所给出的是一种"局部结构指纹"，便于映射回宗族、户籍、八旗等可命名的史学概念；后者基于跨年累积的全局消息传递，所给出的是聚合了远程邻接的"全局信号"，刻画在更大尺度上的人物相似性。两路得分将在第~\ref{chap:mas}~章的多智能体协商协议中分别作为独立证据进入合议，从而形成结构与全局双轨支持的可解释判据。

\section{双半径节点标注}\label{sec:seal-drnl}

封闭子图剥离了原图的全局坐标，因此模型必须从结构本身重新获得"哪些节点在子图中扮演什么角色"的标记。为此，子图链路预测引入双半径节点标注（Double-Radius Node Labeling，DRNL）：对子图内任一节点 $v$，记 $d(v,m)$ 与 $d(v,w)$ 为 $v$ 到两端点 $m,w$ 的无向最短路径长度（取 $\infty$ 表示不可达），令 $d=d(v,m)+d(v,w)$，则节点标号定义为
\begin{equation}\label{eq:drnl}
\ell(v)=
\begin{cases}
0, & d(v,m)=\infty\ \text{或}\ d(v,w)=\infty,\\[3pt]
1, & v=m\ \text{或}\ v=w,\\[3pt]
1+\min\!\big(d(v,m),d(v,w)\big)+\dfrac{d}{2}\!\left(\dfrac{d}{2}+d\bmod 2-1\right), & \text{其他}.
\end{cases}
\end{equation}

式~\eqref{eq:drnl}~具有两点对模型有用的性质。其一，它对 $(d(v,m),d(v,w))$ 严格对称：交换男女两端不会改变标号，从而确保下游学习到的结构表示不偏向任一性别端点。其二，所有具有相同距离对 $(d(v,m),d(v,w))$ 的节点共享同一标号，因而式~\eqref{eq:drnl}~实际上把节点离散化为有限多个"距离 $\times$ 拓扑组合"等价类；后接的关系图卷积网络（Relational GCN，R-GCN）便可以把以"距离 $+$ 拓扑组合"为基的微结构表示作为可学习的嵌入直接吸纳。两端点本身约定 $\ell(m)=\ell(w)=1$，不可达节点 $\ell=0$，使模型能将"目标对"与"同距离的结构邻居"区别对待。

\section{十三类历史微结构图样目录}\label{sec:seal-motifs}

在掩码子图的无向投影上枚举从 $m$ 到 $w$ 的简单路径，并将每条路径上经过的边按可见关系名 $\mathcal{R}_\mathrm{vis}=\{r_\mathrm{hw},r_\mathrm{fs},r_\mathrm{fd},r_\mathrm{sib},r_\mathrm{hh},r_\mathrm{hc},r_\mathrm{cb}\}$ 拼接为关系序列，所得序列称为一个微结构图样（motif）。在中国多代人口面板（CMGPD-LN）训练集 80 个查询事件上，本研究筛选出史学含义明确、识别频率较高的十三类图样，按"父系结构""户与姻亲网络""八旗结构"分为三组并落盘为机读目录（\texttt{seal/motifs/catalog.json}）。完整目录列于表~\ref{tab:motif-catalog}。

\begin{table}[ht]
  \centering
  \caption{十三类历史微结构图样目录（与 \texttt{seal/motifs/catalog.json} 一致；覆盖率为掩码子图上 80 个训练事件中至少存在一条对应路径的比例）}\label{tab:motif-catalog}
  \begin{tabular}{c l l c}
    \toprule
    编号 & 关系序列 & 史学含义 & 覆盖率（\%）\\
    \midrule
    \multicolumn{4}{l}{\emph{父系结构}}\\
    M01 & $(r_\mathrm{sib})$ & 同父登记的兄妹 & 86.7 \\
    M02 & $(r_\mathrm{fs},r_\mathrm{fd})$ & 共同父亲（最关键信号） & 93.3 \\
    M03 & $(r_\mathrm{sib},r_\mathrm{sib})$ & 堂表兄妹 & 80.0 \\
    M04 & $(r_\mathrm{sib},r_\mathrm{sib},r_\mathrm{sib})$ & 大宗族网络 & 80.0 \\
    M05 & $(r_\mathrm{fs},r_\mathrm{fs},r_\mathrm{sib})$ & 父系祖辈兼兄妹（辈分错位） & 80.0 \\
    M06 & $(r_\mathrm{sib},r_\mathrm{fs},r_\mathrm{fd})$ & 共父展开变体 & 80.0 \\
    M07 & $(r_\mathrm{fs},r_\mathrm{fd},r_\mathrm{sib})$ & 共父再延一步兄妹 & 73.3 \\
    M08 & $(r_\mathrm{sib},r_\mathrm{fd},r_\mathrm{fd})$ & 代际错位的同宗联姻链 & 73.3 \\
    \multicolumn{4}{l}{\emph{户与姻亲网络}}\\
    M09 & $(r_\mathrm{hh},r_\mathrm{hh},r_\mathrm{sib})$ & 同户成员之兄妹（包户亲属） & 50.0 \\
    M10 & $(r_\mathrm{hh},r_\mathrm{hh},r_\mathrm{fd})$ & 同户成员之女（主家女儿） & 66.7 \\
    M11 & $(r_\mathrm{sib},r_\mathrm{hw},r_\mathrm{sib})$ & 前代联姻链（重复联姻） & 33.3 \\
    \multicolumn{4}{l}{\emph{八旗结构}}\\
    M12 & $(r_\mathrm{fs},r_\mathrm{cb},r_\mathrm{cb})$ & 八旗内同父系（旗内婚父系视角） & 83.8 \\
    M13 & $(r_\mathrm{cb},r_\mathrm{cb})$ & 同旗（旗内婚） & 85.0 \\
    \bottomrule
  \end{tabular}
\end{table}

第一组 M01--M08 刻画父系登记意义上的同宗结构：M02 是覆盖率最高（93.3\%）的"共同父亲"双边表达，M01 与 M02 在 \texttt{FATHER\_ID} 等价但分别使用派生 $r_\mathrm{sib}$ 与原始 $r_\mathrm{fs},r_\mathrm{fd}$ 表述，因而即便子图内某一类边被屏蔽仍可互补检出。第二组 M09--M11 刻画通过户成员或既有姻亲网络中介的间接联系，其中 M11 的"前代联姻"覆盖率仅 33.3\%，但当其出现时往往对应史学意义上的"门当户对家族"重复联姻，是高价值结构信号。第三组 M12--M13 捕捉八旗制度下的旗内婚习俗，覆盖率均高于 80\%，与清代旗人婚配的史学共识相吻合。

\section{关系图卷积评分模型}\label{sec:seal-rgcn}

掩码子图随后被送入一个四层关系图卷积网络（R-GCN）评分器进行打分。每一节点的输入特征由三段拼接而成：节点类型独热编码（4 维，对应 person/household/community/banner 四类）、双半径节点标号 $\ell(v)$ 的独热编码（$K_\mathrm{DRNL}+1$ 维，含不可达类）以及十六项细粒度数值特征——其中包含目标节点旗标、年龄、存活、性别、累计婚姻数、子嗣数、户规模、社区规模、八旗规模等\citep{fey2019pyg}。这一组合既保留了节点的角色与拓扑等价类信息，也保留了与个人生命史相关的人口学信号，使 R-GCN 在纯结构之外能进一步基于属性差异分辨候选。

第 $\ell+1$ 层的卷积更新满足式~\eqref{eq:rgcn-update}：
\begin{equation}\label{eq:rgcn-update}
h_v^{(\ell+1)}=\sigma\!\left(\sum_{r\in\mathcal{R}_\mathrm{vis}}\sum_{u\in\mathcal{N}_r(v)}\frac{1}{c_{v,r}}\,W_r^{(\ell)}\,h_u^{(\ell)}+W_0^{(\ell)}\,h_v^{(\ell)}\right),\quad W_r^{(\ell)}=\sum_{b=1}^{B}a_{rb}^{(\ell)}V_b^{(\ell)}.
\end{equation}
其中 $\mathcal{N}_r(v)$ 为 $v$ 沿关系 $r$ 的邻居集合，$c_{v,r}=|\mathcal{N}_r(v)|$ 为标准化常数，$\sigma$ 取 ReLU。式~\eqref{eq:rgcn-update}~右端第二项实现关系自环以保留节点自身信息。值得注意的是，关系特定矩阵 $W_r^{(\ell)}$ 经基分解（basis decomposition）由 $B=4$ 个共享基矩阵 $V_b^{(\ell)}$ 线性组合而得；由于基数 $B=4$ 严格小于可见关系数 $|\mathcal{R}_\mathrm{vis}|=7$，不同关系被强制共享部分参数，从而显著缓解了关系数偏多导致的过拟合。

读出阶段采用 SortPooling：将四层输出沿通道维拼接后，按最后隐藏维降序排序并截取前 $k=30$ 个节点的隐藏向量；继而经过一个核宽 2 的 Conv1d 层（输出 16 个通道）与一个核宽 $k-1$ 的 Conv1d 层（输出 32 个通道），最后经两层多层感知机（MLP）映射为标量得分 $s_\theta(m,w,\tau)$。在本研究的配置下，整体模型的可学习参数量共计 97{,}889。

\section{训练与评估}\label{sec:seal-train}

训练采用列表式 softmax 交叉熵损失。对每个真实婚姻事件 $(m,w^\star,\tau)$，从同期合法候选池中按"$10$ 个同社区 $+$ $10$ 个全局"策略采集 $K=20$ 个负样本 $\{w_1^-,\dots,w_K^-\}$，构造候选列表 $\mathcal{Y}=\{w^\star,w_1^-,\dots,w_K^-\}$，最小化
\begin{equation*}
\mathcal{L}(m,w^\star,\tau)=-\log\frac{\exp\!\big(s_\theta(m,w^\star,\tau)\big)}{\sum_{w\in\mathcal{Y}}\exp\!\big(s_\theta(m,w,\tau)\big)}.
\end{equation*}
其中同社区负样本是被刻意采入的"难负样本"，旨在迫使模型不仅仅依赖社区粗粒度信号；优化器为 Adam（学习率 $10^{-3}$，权重衰减 $5\times 10^{-4}$），并以验证集 MRR 连续三轮不上升作为早停准则。

在 80 个 query 的测试集上，模型在 100 候选随机配对池中给出 Hit@1 $=0.975$、MRR $=0.984$ 的评分指标，相较于随机基线 Hit@1 $=0.010$ 提升约 97 倍。该结果与基于全图传播的异构图变换器形成结构层面的相互印证：候选对周围 3 跳之内的局部子图本身已携带极强的婚配判据。

更具诊断意义的是按 $m$ 父系链长度 $\ell_\mathrm{pat}$ 的分组消融。当 $m$ 的登记父亲缺失（$\ell_\mathrm{pat}=0$）时，Hit@1 由父亲已知组的 $1.000$ 急剧下滑至 $0.429$；而仅父亲、祖父及更长链长组合并的群体仍稳定达到 $1.000$。这一显著差异说明：所谓"局部结构指纹"在 CMGPD-LN 上的主导成分仍是父系结构信号——一旦父系链断裂，子图内可被关系图卷积网络利用的判据急剧减少。该发现与第~\ref{chap:hgt}~章母系边消融的结论方向一致，共同支持了本研究关于"详父系、略母系"档案体例下结构信号分布的核心假设。
